\documentclass[12pt,english]{article}
\usepackage[T1]{fontenc}
\usepackage[latin9]{inputenc}
\usepackage[toc,page]{appendix}
\usepackage{geometry}
\geometry{verbose,tmargin=1.5in,bmargin=1.5in,lmargin=1.5in,rmargin=1.5in}
\usepackage{bm}
\usepackage{amstext}
\usepackage{comment}
\usepackage{amsmath}
\usepackage{hyperref}
\hypersetup{
     colorlinks = true,
     citecolor = navy,
     linkcolor = navy,
     urlcolor=navy
}
\usepackage{adjustbox}
\usepackage[space]{grffile}
\usepackage{setspace}
\usepackage[normalem]{ulem}
\usepackage{nicematrix}
\usepackage{xcolor}
\usepackage{soul}
\usepackage{arydshln}
\usepackage{dsfont}
\usepackage{refstyle}
\usepackage[capitalize,nameinlink]{cleveref}
\usepackage{setspace}
\usepackage{bbm}
\usepackage[tableposition=t]{caption} % ensures caption of longtables is normalsized.
\usepackage{mathpazo} % add Palatino font
\usepackage{longtable}
\usepackage{amssymb}
\usepackage{amsfonts}
\usepackage{amsmath}
%\usepackage[margin=1in]{geometry} % to set margins
\usepackage{scalefnt}
\usepackage{wrapfig}
\usepackage{graphicx}
\usepackage{graphics}

%\usepackage{subfigure}
\usepackage{lscape}
\usepackage{pdflscape}
\usepackage{multirow}
\usepackage{array}
\usepackage{rotating}
\usepackage{bibunits}
\usepackage{fix-cm}
\usepackage{float}
\floatplacement{table}{H}
\usepackage{epsfig}
\usepackage{epstopdf}
\usepackage{refstyle}
\usepackage{tikz}
\usepackage{booktabs}
\usepackage{tabularx}
\usepackage{array}
\newcolumntype{x}[1]{>{\raggedleft\arraybackslash}p{#1}}%
\newcolumntype{d}[1]{>{\centering\arraybackslash}p{#1}}%
\newcolumntype{y}[1]{>{\raggedright\arraybackslash}p{#1}}%
\usepackage{verbatim}
%\epstopdfsetup{suffix=}
\setcounter{MaxMatrixCols}{10}
\usepackage[longnamesfirst,authoryear]{natbib}
\newenvironment{proof}[1][Proof]{\noindent \textbf{#1.}}{\rule{0.5em}{0.5em}}
\setlength{\topmargin}{-0.80in} \setlength{\textheight}{9.50in}
\setlength{\evensidemargin}{.0in} \setlength{\oddsidemargin}{.0in}
\setlength{\textwidth}{6.85in}
\setlength{\voffset}{0.4cm}
\setlength{\hoffset}{-0.5cm}
\usepackage{setspace}
\usepackage{afterpage}
\usepackage[flushleft]{threeparttable}
\usepackage{tabulary}
%\usepackage[titletoc,title]{appendix}
\usepackage[bottom]{footmisc}
\usepackage{color}
	\definecolor{maroon}{RGB}{128,0,0}
	\definecolor{navy}{RGB}{0,0,165}
	\definecolor{lb}{RGB}{150,165,190}
	\definecolor{bluenice}{rgb}{0.25,0.51,0.76}
	\definecolor{orangenice}{HTML}{EB811B}
	\definecolor{greennice}{HTML}{14B03D}
\usepackage{colortbl}
\usepackage{caption}
\usepackage{subcaption}
\usepackage{placeins}
\usepackage{nicefrac}
%\usepackage{chngcntr}
\newtheorem{assumption}{Assumption}
\usetikzlibrary{matrix,decorations.pathreplacing}

\newcommand{\sigmac}{\underline{\sigma}_{c,t}}
\newcommand{\sigmad}{\underline{\sigma}_{d,t}}
\newcommand{\dZ}{d\underline{Z}_t}
\newcommand{\sigmacd}{\underline{\sigma}_t^{\Delta}}
\newcommand{\D}{\mathcal{D}_t}
\newcommand{\sigmap}{\underline{\sigma}_{P,t}}
\newcommand{\sigmapc}{\underline{\sigma}_{P_c,t}}
\newcommand{\sigmapd}{\underline{\sigma}_{P_d,t}}
\newcommand{\mup}{\mu_t^{P}}
\newcommand{\mud}{\mu_t^{\mathcal{D}}}
\newcommand{\sigmaD}{\underline{\sigma}_{\mathcal{D},t}}
\newcommand{\muq}{\mu_{q,t}}
\newcommand{\sigmavt}{\underline{\sigma}_{\vartheta,t}}
\newcommand{\sigmaq}{\underline{\sigma}_{q,t}}
\newcommand{\sigmaK}{\underline{\sigma}_{K,t}}
\newcommand{\vt}{\vartheta_t}
\newcommand{\amb}{(\mathds{1}_a\sigma_a-\mathds{1}_b\sigma_b)}
\newcommand{\bma}{(\mathds{1}_b\sigma_b-\mathds{1}_a\sigma_a)}
\newcommand{\nuad}{\underline{\nu}_t^{ad}}
\newcommand{\nubd}{\underline{\nu}_t^{bd}}
\newcommand{\nuab}{\underline{\nu}_t^{a}}
\newcommand{\nubb}{\underline{\nu}_t^{b}}
\newcommand{\sigmaeta}{\underline{\sigma}_{\eta,t}}
\newcommand{\sigmas}{\underline{\sigma}_{S,t}}
\newcommand{\bsigmas}{\underline{\bar{\sigma}}_{S}}
\newcommand{\muvt}{\mu_t^\vartheta}
\newcommand{\PP}{\left(\underline{P}\right)}
\newcommand{\IO}{\mathbf{\Omega\left(\underline{P}\right)}}
\newcommand{\io}{\mathbf{\Omega}}
\newcommand{\fd}{\mathbf{\Xi\left(\underline{P}\right)}}
\newcommand{\fdd}{\mathbf{\Xi}}
\newcommand{\ishare}{\underline{\gamma}\left(\underline{P}\right)}
\newcommand{\ishares}{\underline{\gamma}}
\newcommand{\MM}{\mathbf{\mathcal{M}\PP}}
\newcommand{\MMM}{\mathbf{\mathcal{M}}}
\newcommand{\obj}{\mathbf{\Phi}}
\newcommand{\objt}{\mathbf{\tilde{\Phi}}}
\newcommand{\nuu}{(\underline{\nu})}
\newcommand{\nured}{\textcolor{red}{\nu}}
\newcommand{\nuredOne}{\textcolor{red}{\nu_1}}
\newcommand{\nuredTwo}{\textcolor{orangenice}{\nu_2}}
\newcommand{\scaleBlueOne}{\textcolor{bluenice}{\frac{1-\psi_{PHL,rice\rightarrow 11}\nu}{1-\psi_{PHL,1\rightarrow 11}}}}
\newcommand{\scaleBlueTwo}{\textcolor{bluenice}{\frac{1-\psi_{PHL,rice\rightarrow 1,rice}\nu}{1-\psi_{PHL,rice\rightarrow 1,rice}}}}
\newcommand{\scaleBlueOneOne}{\textcolor{bluenice}{\frac{1-\psi_{PHL,rice\rightarrow 11}\nu_1}{1-\psi_{PHL,1\rightarrow 11}}}}
\newcommand{\scaleBlueTwoOne}{\textcolor{bluenice}{\frac{1-\psi_{PHL,rice\rightarrow 1,rice}\nu_1}{1-\psi_{PHL,rice\rightarrow 1,rice}}}}

\newcommand{\scaleBlueOneTwo}{\textcolor{greennice}{1-\frac{1-\psi_{PHL,rice\rightarrow 11}}{\nu_2 }}}
\newcommand{\scaleBlueTwoTwo}{\textcolor{greennice}{1-\frac{1-\psi_{PHL,rice\rightarrow 1,rice}}{\nu_2}}}
\newcommand{\mb}[1]{\mathbold{\tilde{#1}}_t}
\newcommand{\mbb}[2]{\mathbold{\tilde{#1}}_{#2,t}}

	\definecolor{maroon}{RGB}{128,0,0}
	\definecolor{navy}{RGB}{0,0,165}
	\definecolor{lb}{RGB}{150,165,190}
	\definecolor{MyWhite}{RGB}{255, 255, 255}
	
\definecolor{red}{rgb}{0.9,0.1,0.1}
\definecolor{blue}{rgb}{0,0,0.7}
\definecolor{dgr}{rgb}{0.08,0.25,0.18}
\definecolor{navy}{RGB}{0,0,165}
\definecolor{org}{HTML}{EB811B}
\definecolor{lgr}{HTML}{14B03D}
\definecolor{shadecolor}{rgb}{0.8,0.8,0.8}

\definecolor{graphblue}{RGB}{64,130.56,194.5}
\definecolor{graphred}{RGB}{255,0,0}
\definecolor{graphyellow}{RGB}{237.8340,177.64,32}
\definecolor{graphpurple}{RGB}{126,47,142}

\newcommand{\hlight}[1]{\textcolor{org}{#1}}
\newcommand{\glight}[1]{\textcolor{lgr}{#1}}
\newcommand{\rlight}[1]{\textcolor{red}{#1}}
\newcommand{\blight}[1]{\textcolor{graphblue}{#1}}
\newcommand{\plight}[1]{\textcolor{graphpurple}{#1}}
\begin{document}
\title{CES production -- beyond Cobb-Douglas}
\author{Nick}
\maketitle

\paragraph{Purpose:} The goal of this note is to re-do our calculations of the price of value added and the supply/demand influence calculations when we have CES production instead of Cobb-Douglas.



\section*{Calculating the price of value added}

Production is now given by:

\begin{align*}
	Y_{i,t} &= \left[(1-\gamma_i)VA_{i,t}^{\frac{\epsilon-1}{\epsilon}} + \gamma_i M_{i,t}^{\frac{\epsilon-1}{\epsilon}}\right]^{\frac{\epsilon}{\epsilon-1}}\\
	M_{i,t} &= \left[\sum_j\omega_{j\rightarrow i} M_{j\rightarrow i,t}^{\frac{\sigma-1}{\sigma}}\right]^{\frac{\sigma}{\sigma-1}}
\end{align*}
where $M_{i,t}$ is intermediate inputs used by $i$ and $M_{j\rightarrow i,t}$ are intermediates of product $j$ used by $i$ in production.

Standard optimization gives us the following price index:

\begin{align*}
	P_{i,t} &= \left[(1-\gamma_i)P_{VA,i,t}^{1-\epsilon} + \gamma_i P_{M,i,t}^{1-\epsilon}\right]^{\frac{1}{1-\epsilon}}\\
	P_{M,i,t} &= \left[\sum_j \omega_{j\rightarrow i}P_{j,t}^{1-\sigma}\right]^{\frac{1}{1-\sigma
}}\\
	\Leftrightarrow P_{i,t} &= \left[(1-\gamma_i)P_{VA,i,t}^{1-\epsilon} + \gamma_i \left(\sum_j \omega_{j\rightarrow i}P_{j,t}^{1-\sigma}\right)^{\frac{1-\epsilon}{1-\sigma}}\right]^{\frac{1}{1-\epsilon}}
\end{align*}

Unlike the Cobb-Douglas case, this cannot be linearly inverted!

To solve this we are going to need to do some hat algebra to be able to replace the parameters with data.

So let's do that: I won't derive it but with Hat algebra we can write this in terms of 2017 IO data:

\begin{align*}
	\hat{\underline{P}}_t^{1-\epsilon} &= diag(1-\underline{\gamma}_{2017})\hat{\underline{P}}_{VA,t}^{1-\epsilon} + diag(\underline{\gamma}_{2017}) \left(\mathbf{\Omega}_{2017}\hat{\underline{P}}_{t}^{1-\sigma} \right)^{\frac{1-\epsilon}{1-\sigma}}
\end{align*}
where exponents are elementwise and $\hat{\underline{P}}_t \equiv \underline{P}_{t}/\underline{P}_{2017}$.

Since $\hat{\underline{P}}_t$ is data, calculating $\underline{P}_{VA,t}$ is relatively easy:

\[
	\hat{\underline{P}}_{VA,t} = \left[\hat{\underline{P}}_t^{1-\epsilon} -diag(\underline{\gamma}_{2017}) \left(\mathbf{\Omega}_{2017}\hat{\underline{P}}_{t}^{1-\sigma} \right)^{\frac{1-\epsilon}{1-\sigma}}\right]^{\frac{1}{1-\epsilon}}
\]

Next, we want to think about how we can calculate the influence of different sectors (or collections of sectors). Let's continue to assume that final demand is Cobb-Douglas so that we can write:

\[
	\hat{P}_t^{CPI} = \underline{\omega}_{CPI}' \hat{\underline{P}}_t 
\]

Now we have a problem: this system is non-linear in terms of the underlying prices of value added. Moreover, we don't have a closed form expression for output prices as a function of the price of value added. 

So let me propose the following routine to estimate the contributions of each sector:

\begin{itemize}
	\item We calculate the change in the price of value added with the above formula.
	\item We loop over every sector and time period
	\item We assume that $\hat{P}_{VA,i,t}=1$ (i.e. no change) for just that sector
	\item We use a Newton-Rhapson routine to solve for $P_{CPI,-i,t}$
	\item We then see if these contributions sum to 1 (they won't for a non-linear model). 
	\item The difference from 1 is called interaction
	\item We plot the results.
\end{itemize}



\end{document}